\documentclass[12pt]{article}
\usepackage[utf8]{inputenc}
\usepackage[T1]{fontenc}
\usepackage{hyperref}
\title{Gowli Language Resources for K-12 Curriculum Development}
\author{Compiled by Claude Sonnet 5.5 (claude-sonnet-5-5) for Andr\'as Kornai}
\date{}
\begin{document}
\maketitle

\section{Introduction}
Gowli (ISO 639-3 gok) is, according to Glottolog, a retired code based on a misunderstanding: Gowli appears to be a caste name (Gowari, Gaoli, Gowli, Gowlan) and not a distinct language. Searches found no dictionary, grammar, or teaching material for a language by this name.

\section{Primary Educational Resources}
None found.

\section{Linguistic and Reference Materials}
Glottolog entry discussing the retired code: \url{https://glottolog.org/resource/languoid/iso/gok}. Rosetta/Long Now wiki page: \url{https://rosettapanglossia.longnow.org/index.php?title=Gok}.

\section{Religious and Community Resources}
Joshua Project describes the Gawli, a Yadav sub-caste in northern Karnataka; this may be a different group from the one intended: \url{https://joshuaproject.net/index.php/people_groups/16799}.

\section{Digital and Audio Resources}
None found.

\section{Curriculum Development Context}
No evidence of a separate Gowli language was found, so no native-language curriculum can be based on these sources. Community consultation is needed to determine which language the group actually speaks.
\end{document}